自主反应机理探索需要提出反应步骤、准备三维结构，并确定计算所得鞍点连接哪些反应物和产物。大语言模型（LLM）智能体可能误判空间关系，也可能把计算得到的结构误认为原先想研究的物种或步骤。

本文提出 Cloudra，一个用于反应机理探索的多智能体框架。Delivery DAG 表示研究中的物种与基元步骤。按照“智能克制”的原则，确定性算法匹配反应物和产物中的原子、测量键长和二面角、旋转或平移分子片段，并装配候选过渡态几何。频率计算区分极小值与鞍点，IRC 端点计算确定各鞍点连接的物种。智能体利用这些操作和计算检验候选步骤、比较竞争路径并修订机理。

CloudraBench 用于衡量这套框架的端到端机理计算能力，包含50题，覆盖基元反应、多步机理、立体化学竞争及计算方法局限。四种工作流配置共完成200次运行，Cloudra 的平均综合得分为0.793，同一模型的 Standard 基线为0.581。在 Hard 任务中，平均得分由0.689提高至0.958。40道非 Impossible 任务中，Cloudra 的自洽性评分分量在24题达到满分，Standard 为7题。行为分析进一步考察智能体如何利用键长和二面角测量、优化结构及 IRC 端点修订局部机理判断。
